\subsection{Correcting for description truncation in the transaction-price panel}
\label{sec:truncation}

The realized-transaction panel is assembled from a bulk listings interface that
returns each property's description capped at a fixed length of six hundred and
ninety-nine characters, a limit that truncates between roughly half and two thirds
of descriptions before they are embedded. Because the treatment is the leading
principal component of the description embedding rather than one of the controls,
this truncation perturbs the treatment itself, and left uncorrected it biases the
estimated price effect rather than merely inflating its variance.

We first establish the sign and magnitude of the bias on the full-text corpus,
where the untruncated description is available for the same listings. Recomputing
the treatment from descriptions cut to the same character limit and re-estimating
the effect with the treatment and controls otherwise held fixed, we find that
truncation attenuates the effect toward zero in every market with a detectable
effect, by between fifteen and twenty-seven percent, with the per-market estimates
collected in Table~\ref{tab:truncation}. The distortion is not a byproduct of the
embedding model's own context window, which admits close to four hundred tokens
against a median description of under two hundred, so the character cap discards
text the encoder would otherwise read.

We do not treat this as classical measurement error. The truncation is a
deterministic function of the description that always removes its tail, and the
removed text is longer for more elaborate listings, so the induced error is
correlated with the true treatment and with omitted quality rather than being
independent additive noise. This is the generated-regressor structure of
\citet{pagan1984econometric} and \citet{murphy1985twostep}, and in its modern
data-mined form the problem studied by \citet{yang2020datamined}, rather than the
classical errors-in-variables model. Under non-classical error the direction of the
bias is not guaranteed in advance \citep{bound2001measurement}, and for latent,
text-derived predictors a naive reliability-ratio correction can enlarge the bias
rather than remove it \citep{jerzak2025attenuation}. We therefore measure the
attenuation on validation data instead of assuming it.

The full-text corpus is an external validation sample in the sense of
\citet{chen2008semiparametric} and \citet{carroll2006measurement}, since it carries
both the truncated and the untruncated treatment on the same units. We summarize the
attenuation by the factor $\lambda = \hat\theta_{\text{trunc}}/\hat\theta_{\text{full}}$,
which is scale-free and absorbs the partialling of the controls. Across the six
markets with a detectable effect the factor is stable at $0.79$ with a standard
deviation of $0.05$ and a range from $0.73$ to $0.85$; markets whose effect is
statistically indistinguishable from zero leave the ratio undefined but require no
correction, because a near-null estimate stays near-null under any rescaling and,
as Table~\ref{tab:truncation} shows, none of them is pushed into significance by the
correction. Precision-weighting across all markets, which down-weights the
near-null cases whose ratios are imprecise, gives a pooled factor of $0.81$ with a
standard error of $0.02$, an attenuation of nineteen percent and a correction
factor of $1.24$.

Two checks support applying this factor to the transaction panel, whose treatment
is only available in truncated form. A leave-one-market-out exercise, in which each
market's truncated effect is corrected by the factor estimated on the remaining
markets, recovers the held-out market's untruncated effect to within $0.011$ in the
units of $\theta$, or under six percent, comfortably inside the market-level
confidence intervals; the factor therefore transports across markets. Regression
calibration, which replaces the truncated treatment with its conditional expectation
given the truncated value and the controls and re-estimates the effect, moves each
estimate in the same direction and closes fifty-four percent of the truncation gap.
It under-corrects, which is the expected behavior of regression calibration when the
error is non-classical \citep{carroll2006measurement}, and we accordingly take the
measured attenuation factor as the primary correction and read the calibration
estimate as a conservative floor. As a further robustness option the same paired
validation sample supports a simulation-extrapolation correction
\citep{cook1994simex}, which we do not need here given the stability of the measured
factor.

We report transaction-price effects after dividing the truncated estimate by
$\lambda$, propagating the uncertainty of both the effect and the factor through the
delta method; because the effect and the factor are estimated on disjoint samples
their errors are independent and the corrected standard error is
$\sqrt{\hat\theta_{\text{trunc}}^2\,\widehat{\mathrm{Var}}(\lambda)/\lambda^4
+ \widehat{\mathrm{Var}}(\hat\theta_{\text{trunc}})/\lambda^2}$. To our knowledge
this truncation is not confronted in the earlier real-estate text literature, which
represents descriptions through bag-of-words or dictionary features that carry no
fixed context window \citep{nowak2017textual, shen2021information}; the problem is
specific to fixed-window sentence-embedding features, and the attenuation we
document is the pricing-model counterpart of the classification-accuracy loss that
\citet{sun2019finetune} report for head-only truncation of long documents.

\begin{table}[H]
\centering
\caption{Description-truncation attenuation, measured on the full-text validation
corpus. For each market with a detectable text effect we report the DML effect of
the leading text direction estimated from the untruncated description and from the
description cut to the interface's character limit, both with identical controls,
and the attenuation factor $\lambda=\hat\theta_{\text{trunc}}/\hat\theta_{\text{full}}$.
The pooled factor is precision-weighted across all ten markets and used to correct
the transaction-price panel.}
\label{tab:truncation}
\small
\begin{tabular}{lrrr}
\toprule
Market & $\hat\theta_{\text{full}}$ & $\hat\theta_{\text{trunc}}$ & $\lambda$ \\
\midrule
Boston       & $-0.245$ & $-0.209$ & $0.85$ \\
San Francisco& $-0.202$ & $-0.171$ & $0.85$ \\
Chicago      & $-0.222$ & $-0.175$ & $0.79$ \\
Denver       & $-0.126$ & $-0.099$ & $0.79$ \\
Atlanta      & $-0.138$ & $-0.105$ & $0.76$ \\
Philadelphia & $-0.167$ & $-0.122$ & $0.73$ \\
\midrule
Pooled ($1/\lambda=1.24$) & \multicolumn{3}{r}{$\lambda=0.81\ (0.02)$, attenuation $19\%$} \\
\bottomrule
\end{tabular}
\end{table}
